\chapter*{Agradecimientos}
\addcontentsline{toc}{chapter}{Agradecimientos}

Este proyecto no habría sido posible sin el apoyo y la colaboración de numerosas personas e instituciones que han contribuido, directa o indirectamente, a su realización.

En primer lugar, quiero agradecer a mi tutor de TFG por su orientación académica y su disponibilidad a lo largo de todo el proceso de desarrollo, siendo nuestros horarios tan diferentes. Sus revisiones y conocimiento han sido determinantes para conseguir dar forma a la idea que tenía en mente.

Agradezco también a la empresa Europavia por facilitar el entorno real de despliegue, con una necesidad existente a nivel corporativo que ha servido como banco de pruebas del sistema. La oportunidad de desarrollar un proyecto con impacto práctico inmediato en un contexto profesional ha enriquecido enormemente la experiencia. Sobre todo Nacho y CIVEX por ser siempre de ayuda y facilitar cualquier desarrollo que considerara.

A mis compañeros de carrera, sobre todo a los de Rugby y a aquellos del Erasmus, que me enseñaron el verdadero sentido de la universidad. El camino se hace más llevadero cuando se recorre en buena compañía.

A mi familia, por su paciencia infinita durante los meses más intensos de exámenes. Su apoyo incondicional ha sido el pilar fundamental sobre el que se sustenta todo lo demás.

Por último, quiero reconocer a la comunidad de código abierto, cuyos proyectos --- Ollama, Qdrant, FastAPI, OpenWebUI, Playwright, PaddleOCR, LiteLLM, entre muchos otros --- han hecho posible construir un sistema de esta envergadura sin dependencias comerciales. Este TFG es, en gran medida, un tributo a la filosofía del software libre.

\vspace{1cm}
\begin{flushright}
\textit{Ángel María Ballesteros Limones}\\
\textit{Madrid, 2026}
\end{flushright}
\newpage
